\subsubsection{Subproceso 3.B: Confirmaci\'on y Registro Oficial}

El procedimiento se activa cuando el Subproceso 3.A emite la Pre-Ficha Aprobada o Subsanada: la plataforma bloquea la fase de edici\'on del estudiante y, por la aprobaci\'on autom\'atica del tr\'amite (TUPA PA5400954D, RR-3698-2024), asienta la carga acad\'emica con car\'acter definitivo en la base de datos del SIGA. Desde ese momento el registro se reparte entre dos oficinas: la Oficina de Estad\'istica de cada Facultad conduce la depuraci\'on y el visado; la Direcci\'on de Registro Central y Estad\'istica (DIRCE), dependencia central que custodia la base del SIGA, mantiene el registro oficial. El estudiante ya no dispone de acceso de edici\'on sobre su propia matr\'icula: recibe el Reporte de Matr\'icula como comprobante informativo y, si detecta un error, lo canaliza ante su Facultad.

En este tramo aparece el efecto de los subprocesos previos: los rezagados de reincorporaciones, convalidaciones y rectificaciones tard\'ias llegan a las n\'ominas provisionales cuando la primera semana ya arranc\'o, y hay estudiantes que asisten a c\'atedra sin figurar en los listados.

Al concluir la semana de depuraci\'on, la Oficina de Estad\'istica consolida el padr\'on, entrega a la DIRCE el archivo digital de consolidaci\'on y le comunica que culmin\'o el proceso de matr\'icula de la Facultad (RR-0570-2022, Art. 78\textdegree{}). Con este acto se ejecuta el Cierre de Matr\'icula: ya no se admiten modificaciones, salvo resoluci\'on decanal o rectoral extempor\'anea. El Reporte y la Boleta de Matr\'icula solo tienen efectos legales con el visado institucional de la Oficina de Estad\'istica; si el documento sale observado, la entrega se posterga hasta que se reprocese y vuelva a visar.

En el primer d\'ia h\'abil de la segunda semana, los listados actualizados se ponen a disposici\'on de cada docente para el registro de asistencias y calificaciones (RR-0570-2022, Art. 78\textdegree{}). La DIRCE consolida el padr\'on nacional de matriculados y remite los reportes a SUNEDU y MINEDU dentro del plazo legal de veinte d\'ias (RR-0570-2022, Art. 79\textdegree{}); tambi\'en habilita el tr\'amite de expedici\'on del carn\'e universitario. En un plazo de hasta diez d\'ias h\'abiles tras el cierre, las Facultades reciben el registro de matr\'icula, las boletas y la plataforma del SIGA actualizada (RR-0570-2022, Art. 80\textdegree{}). Con ese conjunto de salidas la matr\'icula del periodo queda oficializada.

Los diagramas muestran dos problemas documentados en el registro de campo del periodo 2026-2. El primero: el cierre se comunica solo a t\'itulo informativo; la DIRCE lo acepta sin contrastar el padr\'on remitido con la n\'omina asentada. El segundo, la deficiencia m\'as reportada del levantamiento: la Facultad no puede editar la matr\'icula ya asentada. Los ajustes de horario acordados con alumnos y docentes quedan fuera de su alcance, y cada correcci\'on exige una solicitud manual a la DIRCE con el reproceso y reemisi\'on correspondiente.

Esta descripci\'on del proceso actual sirve de base para el diagn\'ostico y la propuesta de mejora que se desarrollan en las siguientes secciones del informe.
